\section{结论与展望}

\subsection{研究结论}

本文围绕"多模态图检索增强的智能体工作流自主生成"这一核心命题，构建了从多模态知识图谱构建到可执行工作流输出的完整方法体系，并通过 30 条电力领域基准任务的全量实验对各模块进行了独立量化分析。研究结论总结如下：

（1）\textbf{多模态知识图谱构建方法（算法 A）。}提出面向电力领域的多模态文档解析流水线，支持文本、图像和流程图三种模态的统一实体-关系抽取，利用 BGE-M3 嵌入模型（1024 维）生成语义向量，基于 Neo4j 图数据库建立 HNSW 向量索引，最终构建包含 4{,}337 个节点和 5{,}228 条关系的电力知识图谱。该图谱为后续检索与匹配提供了一致的领域知识底座。当前实验中图像与流程图链路尚未在端到端任务上独立验证，相关多模态收益的量化将在未来工作中开展。

（2）\textbf{基于子图扩展的 GraphRAG 检索方法（算法 B）。}提出"语义向量检索种子节点 + BFS 子图扩展 + LLM 重排序"的两阶段检索策略，并设计结构化的图-文上下文格式用于 LLM 推理。实验结果显示：在当前 4{,}337 节点的图谱规模下，向量种子检索（vector\_rag）已能为下游匹配提供足够上下文，BFS 子图扩展（$h=2$）反而平均下降 12.9 个百分点的角色分配准确率，并伴随约 7.4~s 的额外延迟。案例分析表明，扩展引入的次要相关节点会改变 LLM 对任务结构的理解，从而扰动子任务粒度。该结果提示子图扩展的有效场景需要配合自适应跳数控制和更大规模图谱，相关方法本研究保留为未来研究方向。

（3）\textbf{子任务-智能体混合匹配算法（算法 C）。}提出三维加权评分函数，从能力嵌入相似度（$\alpha=0.5$）、工具覆盖率（$\beta=0.3$）和输入输出兼容性（$\gamma=0.2$）三个维度系统化匹配子任务与智能体，并设计软约束回退机制保证全覆盖。实验表明，算法 C 是 RAA 的决定性因素：在保留 BFS 扩展的条件下，算法 C 将 RAA 从 0.056（随机分配基线）提升至 0.296（+24.0 pp）；在不使用 BFS 扩展的条件下，提升幅度进一步扩大至 +42.5 pp，使 RAA 达到 0.425。该模块在 OPS、KQA、DA 三类场景下均稳定贡献准确率提升，是本研究的主要技术贡献。

（4）\textbf{工作流 DAG 编排与验证机制（算法 D）。}设计基于 LLM 的隐式数据流依赖推断方法构建工作流 DAG，并实现 DAG 无环性、接口兼容性和任务完整性三层迭代验证-修复机制（最大迭代次数 5）。实验中所有方法的工作流可执行率均达到 100\%，验证了 DAG 验证-修复闭环的有效性。

综上，本研究证明了"基于图谱的语义检索 + 三维量化智能体匹配"是提升专业领域工作流自动生成质量的有效路径，特别是算法 C 在多个场景下表现出稳定且可解释的角色分配能力，为电力调度、故障应急处置、知识服务等场景提供了可行的工程参考。

\subsection{研究局限性}

本研究在实验环节坚持以全量数据如实呈现各模块的真实效果，揭示出以下局限性：

\textbf{（1）BFS 子图扩展未带来稳定增益。}当前 4{,}337 节点的图谱规模相对有限，BFS $h=2$ 扩展引入的节点中相当一部分与查询主题弱相关，反而稀释了下游匹配信号。这一观察提示子图扩展并非"越多越好"，需要结合查询置信度自适应控制跳数，并在更大规模、更稀疏的图谱上重新评估其增益边界。

\textbf{（2）LLM-as-Judge 存在系统性偏差。}Logic 指标采用 LLM 裁判时，简洁、表述通顺的工作流被一致性地高估（pure\_llm 取得 0.947 的 Logic 分但 RAA 为 0.000），与 RAA 排序完全相反。本研究在结论中将 RAA 与 SC 视为更可靠的客观指标。Logic 指标的修正应通过引入人工评分、多模型集成裁判或细化裁判提示词等方式实现。

\textbf{（3）多模态链路未在端到端任务上独立量化。}算法 A 的图像与流程图处理路径已设计完毕，但当前实验主要使用文本来源的 PowerGridQA 与 OPSD 构建图谱节点。多模态信息在工作流生成任务上的独立增益（例如增加单线图后 OPS 任务 RAA 是否提升）尚未单独验证，是评估的一个空白。

\textbf{（4）评估规模与领域覆盖。}本文已将初始 6 条人工标注任务扩展为 30 条分层基准任务，但相较真实电网运行中的故障类型、知识问题和数据分析需求，样本规模仍然有限。未来研究应继续扩展至数百条任务，并引入自动化预标注与人工复核流程降低标注成本。

\textbf{（5）实时性与资源消耗。}算法 C 的 IO 兼容性维度当前以静态字符串匹配实现，相较论文方法形式中的 LLM 判别版本牺牲了一些精度但大幅降低了 LLM 调用频次；算法 D 的依赖推断和验证修复仍需 LLM 调用，端到端延迟约为 9 至 22 秒。对于实时性要求高的场景（如电网事故紧急处置），需进一步通过缓存与本地小模型替换降低延迟。

\subsection{未来研究方向}

基于上述局限性，未来研究可从以下方向展开：

\textbf{（1）自适应子图扩展策略。}本研究观测到固定 $h=2$ 的 BFS 扩展未能稳定提升下游质量。未来可设计基于查询置信度的自适应跳数选择器（$h\in\{0,1,2,3\}$）、基于关系类型的方向性遍历，或将 GraphRAG 中的社区摘要机制与子图扩展结合，让子图规模随任务复杂度动态调整。

\textbf{（2）多模态链路独立验证。}补充包含电力单线图、保护配置图和操作流程图的图像数据集，量化多模态节点对任务分解和角色匹配的独立贡献，验证算法 A 中视觉语言模型 + 实体抽取 + 跨模态对齐流水线的端到端价值。

\textbf{（3）IO 兼容性的语义化重构。}将算法 C 中的字符串包含式 IO 兼容性升级为基于嵌入的语义判别或 LLM 三级评分，并通过缓存机制控制 LLM 调用次数，在保持精度的同时降低延迟。

\textbf{（4）工作流执行反馈闭环。}引入强化学习或经验回放机制，将工作流实际执行结果（步骤成功/失败、输出质量）反馈给匹配算法，实现 $\alpha,\beta,\gamma$ 权重参数的自适应优化。

\textbf{（5）跨领域通用性验证。}将所提方法迁移至医疗诊断、法律文书分析和金融风控等其他专业领域，验证三维量化匹配在不同图谱结构下的稳定性，并研究领域适配所需的最小知识图谱规模。

\textbf{（6）端到端工作流执行验证。}当前实验仅评估工作流生成质量，未对 LangGraph 导出的工作流进行实际执行验证。未来可搭建电力系统仿真环境，接入真实智能体并执行生成工作流，以任务最终完成率为终极评估指标，形成更闭环的评估体系。
